% --- Patron: The Traveler (Ways & Journeys) ---
\subsection{The Traveler (Ways \& Journeys)}

\textit{Lore.} The Traveler is the eternal guide of the road, guardian of those who walk the paths between what is and what might be. Among the Fhara caravans and Kuvani traders along the Way of Silk, the Traveler is invoked at every crossroads, honoured with small offerings at each waypoint, and consulted before every major journey.

The Traveler is not merely one who shows the way---they \emph{are} the way, existing in the pause between steps and in the choice of which path to take when roads fork. Every journey is both physical and spiritual; to move from one place to another is to transform, and the road itself becomes a teacher.

Across the known world, the Traveler wears many faces. In Kahfagia, the Traveler is the \textit{Lantern‑Bearer}, whose light guides ships through uncharted waters. In Ykrul, the Traveler is the \textit{Hoof‑Print}, the spirit of the steppe road that appears as a single arrowhead left at a crossroads. In Aeler, the Traveler is the \textit{Ways‑Stone}, carved into every mile marker by True‑Masons to ensure the road remembers those who walk it. In the Mistlands, the Traveler is invoked with a single unknotted bell‑rope -- a sign that the fog may part for one who walks humbly. Wherever a traveller sets foot, the Traveler is already there, waiting to see which road they choose.

\begin{quote}
One foot in a promise, and the road will meet you halfway. But break your word to the way, and the way will break you.
\end{quote}

\noindent\textbf{Domain Focus}
\begin{itemize}
  \item \textbf{Wayfinding:} navigation, route‑finding, safe passage
  \item \textbf{Journeys:} travel, pilgrimage, exploration
  \item \textbf{Thresholds of Movement:} doors, bridges, crossroads, borders
  \item \textbf{Road Covenants:} promises made on the road, hospitality to travellers
\end{itemize}

\subsubsection*{Thiasos and Codex (Runekeeper Options)}
A Runekeeper sworn to the Traveler binds a creature of the open road. Their \textbf{Thiasos} is a small beast or mythological being that embodies movement and guidance: a grey hare that never tires, a sparrow that knows every shortcut, a small fox that leaves no tracks, a dung beetle that rolls its ball unerringly toward the nearest water, or a tiny brass compass that hums when the bearer is on the true path. This creature requires no food, but it must walk at least a mile each day (or be allowed to roam freely in a new place). If confined for more than a day, it becomes \textit{Compromised} (it grows lethargic and its guidance becomes unreliable, imposing –1 die on all Navigation and Survival rolls for its bearer until it is exercised).

The \textbf{Codex} of the Traveler is never a stationary book. It may be a leather strap knotted with directions, a set of scrolls kept in different saddlebags, a map that redraws itself after each journey, or a pilgrim's diary whose blank pages fill with route descriptions as the bearer walks. Because the road resists being captured in ink, the Codex requires \textit{upkeep}: the Runekeeper must spend an hour each week walking a new path and annotating it, or the Rites stored there become faded (treat as Compromised until updated).

\subsubsection*{Patron's Gift (Imbuement)}
Once per scene as an action, a Runekeeper of the Traveler may touch a held item (weapon, tool, or cloak) and whisper a waymark. For the rest of the scene, the item grants \textbf{+1 die} to \textit{Navigation} or \textit{Survival} rolls, and the first time each scene the bearer would be slowed by difficult terrain or a bureaucratic delay, they may ignore it (the road opens for them). After the scene, the user marks \textbf{+1 Obligation} as the road’s claim on them deepens.

\subsubsection*{Rites of the Traveler}

\paragraph{Rite of Road-Sense (Low, 4 XP)} \hfill \\
\emph{Scene; Self; No.} \\
\textbf{Tags:} \texttt{[NAVIGATION]}, \texttt{[AWARENESS]} \\
\textbf{Materials:} A road‑nail or waystone pebble. \\
\textbf{Effect:} Unerringly pick the fastest safe route in Near/Far. Gain \textbf{+1 die} to avoid ambushes or delays. Create a 4‑segment \emph{Path Memory} clock to ignore difficult terrain once. \\
\textbf{Push It:} Spot one hidden bypass, but generate \textbf{1 SB (Clubs)}. \\
\textbf{Invoke:} 1 action; mark \textbf{+1 Obligation}. \\
\emph{Requires: Familiar (\textit{Invoke:} 1 Boon).}

\paragraph{Rite of the Traveler's Boon (Low, 5 XP)} \hfill \\
\emph{Scene; Self/Ally; No.} \\
\textbf{Tags:} \texttt{[MOVEMENT]}, \texttt{[BOND]}, \texttt{[HASTE]} \\
\textbf{Materials:} Thread tied at the wrist. \\
\textbf{Effect:} Ignore one level of difficult terrain or bureaucracy; \textbf{+1 Effect} to travel or escape checks. If shared, create a 2‑segment \emph{Shared Journey} bond with an ally. \\
\textbf{Push It:} Extend to one more ally; mark \textbf{1 SB (Diamonds)}. \\
\textbf{Invoke:} 1 action; mark \textbf{+1 Obligation}. \\
\emph{Requires: Familiar (\textit{Invoke:} 1 Boon).}

\paragraph{Rite of the Waymark (Standard, 8 XP)} \hfill \\
\emph{Scene; Near; No.} \\
\textbf{Tags:} \texttt{[MARK]}, \texttt{[GUIDANCE]}, \texttt{[PATH]} \\
\textbf{Materials:} A chalk mark or small cairn. \\
\textbf{Effect:} Declare a lane as permitted or easy. Allies gain improved Position/Effect or ignore one obstacle. Create a 6‑segment \emph{Marked Path} clock. \\
\textbf{Push It:} The lane persists between scenes; the first enemy to exploit it forces \textbf{1 SB (Spades)}. \\
\textbf{Invoke:} 1 action; mark \textbf{+1 Obligation}. \\
\emph{Requires: Familiar + Codex (\textit{Invoke:} 1 Boon).}

\paragraph{Rite of the Bridge Between (Standard, 7 XP)} \hfill \\
\emph{Instant; Near; No.} \\
\textbf{Tags:} \texttt{[TELEPORT]}, \texttt{[PATH]}, \texttt{[BOND]} \\
\textbf{Materials:} Two pinches of road‑dust clapped. \\
\textbf{Effect:} Relocate a willing target within Far along a visible or named route. Unwilling targets may resist (Body+Resolve DV 3). Create a 4‑segment \emph{Pathway Established} clock. \\
\textbf{Push It:} Carry one extra ally or bundle; arrivals are off‑balance. \\
\textbf{Invoke:} 1 action; mark \textbf{+1 Obligation}. \\
\emph{Requires: Familiar + Codex (\textit{Invoke:} 1 Boon).}

\paragraph{Rite of the Crown of Crossings (High, 13 XP)} \hfill \\
\emph{Scene; Zone; No.} \\
\textbf{Tags:} \texttt{[AREA]}, \texttt{[COMMAND]}, \texttt{[MOVEMENT]}, \texttt{[FEAR]} \\
\textbf{Materials:} A brass compass missing its needle. \\
\textbf{Effect:} Call the Road: allies gain \textbf{+1 die} to move or evade; pursuers suffer \textbf{--1 die}. Once, declare "the long way is short" to finish a travel clock segment for free. Enemies must test (Wits+Command DV 4) or suffer \textbf{--1 die} to movement. \\
\textbf{Push It:} Seal a hostile route briefly; generate \textbf{2 SB (Clubs/Diamonds)}. \\
\textbf{Invoke:} 1 action; mark \textbf{+2 Obligation}. \\
\emph{Requires: Familiar + Codex + Tier III (\textit{Invoke:} 2 Boons).} \\
\emph{Obligation:} 7 segments.

\paragraph{Rite of the Wayfarer's Covenant (High, 14 XP)} \hfill \\
\emph{Extended; Near; No.} \\
\textbf{Tags:} \texttt{[VOW]}, \texttt{[COMMAND]}, \texttt{[BOND]}, \texttt{[FEAR]} \\
\textbf{Materials:} Waystones from multiple regions. \\
\textbf{Effect:} Bind present parties to safe passage. While honoured: \textbf{+2 Effect} on travel, reroll one failed travel roll per scene. Breaking the oath inflicts Harm 1 (Fatigue) and marks the breaker as \emph{Oathbreaker of the Road} (--2 dice to travel rolls). \\
\textbf{Push It:} Breaking the covenant inflicts Harm 2 (Stress) and attracts hostile Wayward spirits. \\
\textbf{Invoke:} Extended rite; mark \textbf{+2 Obligation}. \\
\emph{Requires: Familiar + Codex + Tier III (\textit{Invoke:} 2 Boons).} \\
\emph{Obligation:} 7 segments.

\subsubsection*{Cantors \& Cults of the Traveler}

\textbf{The Cantor's Song.} A Cantor of the Traveler sings in \textit{verses of the open road} – melodies that rise and fall like hills, harmonies that echo the creak of wagon wheels. Their voice is used to mark a crossroads with a sung sign, to calm a spooked pack animal, or to call a lost wanderer back to the path. The Traveler's Cantors are often found in caravan camps, at pilgrimage shrines, and walking the roads alone, humming to keep the miles company.

\textbf{The Cult of the Unbroken Way.} The Traveler's cults gather at forks, at bridges, and at the crests of passes – never in buildings. The Unbroken Way is a fellowship of guides, couriers, scouts, and those who have walked so far that the road has become their home. A ritual of the Cult involves each member placing a small stone at a crossroads, then walking a different way home than they came – the road must be honoured with new steps. Their creed is scratched into a single waystone that travels from meeting to meeting: \textit{"The road does not end. It only turns."} Outsiders are welcome to walk a mile with the cult, but they must carry a stone from home and leave it at the next crossroads.

\subsubsection*{Witchcraft of the Traveler (Craft / Wayfinding)}

A witch who walks with the Traveler is a \textit{map‑scribe} or \textit{waystone‑carver} – one who crafts the physical markers of the road: a compass that always points to a chosen destination, a map that shows routes that appear only when needed, or a staff that never causes the bearer to trip. Their magic works through the materials of travel: a quill that draws paths that become real by morning, a chisel that carves a symbol into a stone so that travellers who touch it feel rested, or a needle that sews a patch into a cloak that keeps the wind from shifting direction. In Aeler, a True‑Mason might carve a "way keystone" into a bridge so that the bridge remembers every traveller and wishes them well. In Vhasia, a cartographer might illuminate a map with a prayer; the map will never show the route to an ambush.

\textbf{The Witch's Tool: The Way Chisel.} A chisel of iron, its head struck with a simple arrow. Once per session, the witch may carve a "way mark" into a stone, a tree, or a wooden post. For the next day, anyone who sees the way mark gains a perfect sense of direction to the nearest safe harbour, village, or crossroads – but only while walking. The chisel can also be used to "erase" a false trail, removing tracks from a path as if they had never been made (one exchange of effort per use).

\textbf{A Hedge Gift: Road‑Worn Token (2 XP).} Once per scene, the witch may touch a small stone, a coin, or a button that has been carried on a long journey. For the next hour, the token glows faintly when the bearer is on the correct path to their goal (as long as the goal is a physical location). The token does not reveal the path – only that the bearer has not (or has) strayed from it.

\subsubsection*{The Traveler's Corruption Table}

\begin{tblr}{
  colspec = {Q[c,1.5cm] X[2.5] X[2.5]},
  rowhead = 1,
  row{odd} = {bg=gray!10},
  row{1} = {bg=gray!30, font=\bfseries},
  hlines,
}
Tier & Benefit & Cost / Quirk \\
1 & \textbf{Pathfinding Instinct:} +1 die to Navigation and movement rolls in familiar territory. & \textbf{Restless Spirit:} Suffer 1 Fatigue when forced to remain stationary or in confined spaces for extended periods. \\
2 & \textbf{Wayfarer's Luck:} Once/scene, treat a failed travel or navigation roll as a success, but mark 1 SB (Clubs). & \textbf{Compulsive Wandering:} Must take the less obvious path when choices present themselves; suffer --1 die when forced to follow direct routes. \\
3 & \textbf{Road Mastery:} +2 dice to rolls involving travel, escape, or finding hidden paths. & \textbf{Travel Dependency:} --1 die to rolls requiring local knowledge or settlement‑based skills. \\
4 & \textbf{Journey's Blessing:} Once/session, grant all allies within Near range +1 die to travel‑related rolls for one scene. & \textbf{Nomad's Isolation:} --1 die to social rolls involving long‑term relationships or community ties. \\
5 & \textbf{Waypoint Sense:} Once/session, instantly know the location and nature of all significant paths, routes, or crossroads within a day's travel. & \textbf{Directional Obsession:} --1 die to rolls not involving movement, travel, or pathfinding; always aware of cardinal directions. \\
6+ & \textbf{Master of All Ways:} Once/session, become one with the road network. For one scene, teleport between any known waypoints and grant +2 dice to all travel‑related actions. & \textbf{Eternal Journey:} Mark +3 Obligation when using this power; risk permanent Harm (Stress) from the endless call of the road, making settlement impossible until the journey's end. \\
\end{tblr}

\subsection*{Playstyle Notes}
The Traveler is the ideal patron for beginners and veterans alike – its rites are forgiving, its Obligation low, and its benefits immediately useful in almost any campaign. Followers become masters of navigation, escape, and the subtle art of being in the right place at the right time. The Traveler does not demand worship; it asks only that travellers respect the road and keep their bargains. Intrusions often take the form of a sudden fork in the path, a helpful stranger who knows too much, or a road that leads somewhere unexpected. Ideal for guides, couriers, scouts, pilgrims, and anyone who has ever felt the pull of the horizon.

\noindent\textbf{Emphasises}
\begin{itemize}
  \item \textbf{Wayfinding:} the art of never being lost
  \item \textbf{Road Covenants:} a promise made on the road is binding
  \item \textbf{Journey as Transformation:} the road changes the walker
  \item \textbf{Thresholds of Movement:} doors, bridges, borders, and the spaces between
  \item \textbf{Hospitality to Travellers:} the road gives, and the road expects
\end{itemize}